\scanpage{96}
éléments $2_\Lambda$ et $0_\Lambda=4_\Lambda$. Notons que comme
$a^{-1}$ inv. et $2\ne0$, on a $c=a^{-1}\cdot2\ne0$ d'où $a^{-1}2=2$,
donc
\[
c=2
\]

Donc on a deux cas

1) Si \deleted{$\Lambda$ \unsure{est un anneau|est une} $4_\Lambda=0$, $2_\Lambda\ne0$, et $2\cdot\Lambda$ \unsure{réduit|est}} l'idéal
$2\Lambda$ n'a que deux éléments $0$ et $2$ (i.e.
$\Lambda$ est une extension d'une $\mathbb F_2$-algèbre \inserted{$\Lambda_0$} par
un idéal engendré par $\unsure{2\cdot1|2_\Lambda|2}$ et ayant $2$ comme
seul élément non nul) alors
% slanted
\margin{Ce cas a lieu où $\Lambda$ est l'image inverse de l'extension $\mathbb Z/4\mathbb Z$ de $\mathbb F_2$ par $\mathbb F_2$, par un hom de $\mathbb F_2$-alg. $\Lambda_0\to\mathbb F_2$}
\[
N_{G'}(H')=\left\{p\begin{pmatrix}a&b\\c&d\end{pmatrix}\ \middle|\ 
\begin{gathered}\deleted{\unsure{a\ne0}},ad-bc=1\\c\in\{0,2\}\end{gathered}\right\}\tag*{(54)}
\]
admet $N^0_{G'}(H')\overset{\text{déf}}{=}\left\{p\left(\begin{smallmatrix}\lambda&b\\0&\lambda^{-1}\end{smallmatrix}\right)\mid\lambda\in\Lambda^*,b\in\Lambda\right\}$
comme sous-groupe d'indice $2$.
% slanted
\margin{en fait, si $\sigma=p\left(\begin{smallmatrix}1&0\\2&1\end{smallmatrix}\right)$ alors $N_{G'}(H')$ est produit \inserted{$1/2$ direct} de $N^0_{G'}(H')$ par $\{1,\sigma\}$. [Le $\sigma$ ou antipodisme a un sens intrinsèque, indépend$^{\text{mt}}$ du choix d'une base}

2) Dans tout autre cas, \deleted{(p. ex. si $4_\Lambda\ne0$),}
\[
N_{G'}(H')=N^0_{G'}(H')\tag*{(55)}
\]

Dans ce cas \inserted{$2^\circ$}, on voit donc que le commutant
de $\mathrm{Sl}_\Lambda(M)'$ dans $\unsure{S|\widetilde S|\mathfrak S}_M=M^*/\pm1$ est $Z'_\Lambda$.
Dans le cas $1^\circ$) par contre, $Z'_\Lambda$ est d'indice
$2$ dans ce commutant, qui est isom. au
produit direct \inserted{(pas seulement $1/2$ direct!)} de $Z'_\Lambda$ par le sous-groupe
$\{1,\sigma\}$. Le cas typique est celui de
$\Lambda=\mathbb Z/4\mathbb Z$, \deleted{$\sigma(\dot x)=(x+2y)^{\displaystyle\cdot}={}$} alors $Z'_\Lambda=\{1\}$,
mais le groupe des autom de $S_M$ (ens des
sommets de l'octaèdre $P_M$) qui commutent à
$\operatorname{Aut}^{+}(P_M)$, n'est pas réduit à $\{1\}$, il est engendré
par l'antipodisme de l'octaèdre.
% archive sheet number: 95
